\documentclass[12pt,a4paper]{article}

% --- Packages ---
\usepackage[margin=1in]{geometry}
\usepackage{graphicx}
\usepackage{booktabs}
\usepackage{listings}
\usepackage{xcolor}
\usepackage{hyperref}
\usepackage{amsmath}
\usepackage{float}

% --- Code listing style ---
\lstset{
  basicstyle=\ttfamily\small,
  breaklines=true,
  frame=single,
  backgroundcolor=\color{gray!10},
  keywordstyle=\color{blue},
  commentstyle=\color{green!50!black},
  stringstyle=\color{red},
  showstringspaces=false
}

% --- Title ---
\title{
  \vspace{-1cm}
  \textbf{DAS 839 -- NOSQL SYSTEMS} \\[0.5em]
  \Large Assignment -- II \\[0.3em]
  \large MapReduce \& Apache Hadoop
}
\author{
  Mayank Tanwar (IMT2023012) \\
  Areen Patil (IMT2023013) \\
  Kunal Jindal (IMT2023049) \\
  Ansh Gupta (IMT2023540)
}
\date{}

\begin{document}
\maketitle

% ====================================================================
% PROBLEM 1
% ====================================================================
\section{Problem 1: Co-Occurring Word Matrix Generation}

\subsection{Problem 1(a): Top 50 Most Frequent Words (3 Marks)}

\subsubsection{Program Logic}
We implemented a \textbf{two-job MapReduce chain} in the class \texttt{Top50Words.java}.

\textbf{Job 1 -- Word Frequency Count:}
\begin{itemize}
  \item The \texttt{TokenizerMapper} loads the \texttt{stopwords.txt} file via Hadoop's \textbf{Distributed Cache} mechanism during the \texttt{setup()} phase.
  \item Each line of input text is tokenized by splitting on non-alphabetic characters (\texttt{[{\textasciicircum}a-z]+}).
  \item Tokens matching any stop-word are filtered out. Valid tokens are emitted as \texttt{(word, 1)} pairs.
  \item The \texttt{IntSumReducer} aggregates counts for each word.
  \item A \texttt{Combiner} (same as reducer) is also attached at the map-class level for local optimization.
\end{itemize}

\textbf{Job 2 -- Top 50 Extraction:}
\begin{itemize}
  \item A Java \texttt{PriorityQueue} (min-heap of size 50) is maintained locally in each mapper's memory.
  \item Each mapper emits its local top-50 candidates to a \textbf{single reducer} (\texttt{numReduceTasks = 1}).
  \item The reducer applies the same min-heap logic globally to extract the final top 50 words.
  \item Output is written in descending frequency order.
\end{itemize}

\subsubsection{Pseudocode}
\begin{lstlisting}
// Job 1: Mapper
setup(): load stopwords from DistributedCache into HashSet
map(key, line):
    for each token in line.split("[^a-z]+"):
        if token not in stopWords:
            emit(token, 1)

// Job 1: Reducer
reduce(word, counts):
    emit(word, sum(counts))

// Job 2: Mapper
setup(): minHeap = new PriorityQueue(size=50)
map(word, count):
    minHeap.add((word, count))
    if minHeap.size > 50: minHeap.poll()
cleanup():
    for each (word, count) in minHeap: emit(NULL, "word\tcount")

// Job 2: Reducer (single reducer)
reduce(NULL, values):
    globalHeap = new PriorityQueue(size=50)
    for each value: globalHeap.add(parse(value))
    output globalHeap in descending order
\end{lstlisting}

\subsubsection{Output}
The top 50 most frequently occurring non-stop-words are shown below (truncated):

\begin{lstlisting}
quot    1537        nintendo  338
apos     954        http      314
s        880        this      255
lt       503        category  240
gt       495        www       221
\end{lstlisting}

\subsubsection{Screenshots}
\begin{figure}[H]
  \centering
  \includegraphics[width=0.9\textwidth]{screenshots/1a_output_1.png}
  \caption{Problem 1(a): Top 50 Words Output -- Top portion}
\end{figure}

\begin{figure}[H]
  \centering
  \includegraphics[width=0.9\textwidth]{screenshots/1a_output_2.png}
  \caption{Problem 1(a): Top 50 Words Output -- Bottom portion}
\end{figure}

% ---------------------------------------------------------------
\subsection{Problem 1(b): Co-Occurrence Matrix -- Pairs Approach (3 Marks)}

\subsubsection{Program Logic}
Implemented in \texttt{CoOccurrencePairs.java}.

\begin{itemize}
  \item The mapper loads the Top 50 word list via Distributed Cache into a \texttt{HashSet} during \texttt{setup()}.
  \item For each token at position $i$ in the input line, we examine tokens at positions $i+1$ through $\min(i+d, n-1)$, where $d$ is the configurable distance parameter.
  \item Only pairs where \emph{both} words belong to the Top 50 list are emitted.
  \item Pairs are canonicalized by lexicographic ordering to avoid duplicates: if $u < v$, emit $(u,v)$; otherwise emit $(v,u)$.
  \item The reducer sums all counts for each unique pair.
\end{itemize}

\subsubsection{Pseudocode}
\begin{lstlisting}
setup(): load top50 words into HashSet; d = config.distance
map(key, line):
    tokens = line.split("[^a-z]+")
    for i = 0 to len(tokens):
        u = tokens[i]
        if u not in top50: continue
        for j = i+1 to min(i+d, len-1):
            v = tokens[j]
            if v not in top50 or u == v: continue
            pair = (min(u,v), max(u,v))
            emit(pair, 1)

reduce(pair, counts):
    emit(pair, sum(counts))
\end{lstlisting}

\subsubsection{Runtime Analysis}
\begin{table}[H]
  \centering
  \begin{tabular}{@{}lcccc@{}}
    \toprule
    \textbf{Approach} & $d=1$ & $d=2$ & $d=3$ & $d=4$ \\
    \midrule
    Pairs (baseline) & 1537 ms & 1546 ms & 1554 ms & 1544 ms \\ 
    \bottomrule
  \end{tabular}
  \caption{Pairs approach runtime for varying word distance $d$}
\end{table}

\subsubsection{Screenshots}
\begin{figure}[H]
  \centering
  \includegraphics[width=0.9\textwidth]{screenshots/1b_run_1.png}
  \caption{Problem 1(b): Pairs approach -- MapReduce execution start}
\end{figure}

\begin{figure}[H]
  \centering
  \includegraphics[width=0.9\textwidth]{screenshots/1b_run_2.png}
  \caption{Problem 1(b): Pairs approach -- Reduce phase and runtime report}
\end{figure}

\begin{figure}[H]
  \centering
  \includegraphics[width=0.9\textwidth]{screenshots/1b_output.png}
  \caption{Problem 1(b): Pairs co-occurrence matrix output sample}
\end{figure}

% ---------------------------------------------------------------
\subsection{Problem 1(c): Co-Occurrence Matrix -- Stripes Approach (3 Marks)}

\subsubsection{Program Logic}
Implemented in \texttt{CoOccurrenceStripes.java}.

\begin{itemize}
  \item Instead of emitting individual pairs, the mapper builds an \textbf{associative array} (\texttt{MapWritable}) for each word $u$.
  \item The stripe contains all neighboring Top 50 words within distance $d$ (both directions: $i-d$ to $i+d$) as keys, with their co-occurrence counts as values.
  \item This dramatically reduces the number of key-value pairs emitted compared to the Pairs approach.
  \item The reducer merges all stripes for each word by element-wise addition of counts.
\end{itemize}

\subsubsection{Pseudocode}
\begin{lstlisting}
map(key, line):
    tokens = line.split("[^a-z]+")
    for i = 0 to len(tokens):
        u = tokens[i]
        if u not in top50: continue
        stripe = {}
        for j = max(0, i-d) to min(len-1, i+d):
            if j == i: continue
            v = tokens[j]
            if v not in top50 or u == v: continue
            stripe[v] += 1
        if stripe not empty: emit(u, stripe)

reduce(word, stripes):
    merged = {}
    for each stripe: merged += stripe  // element-wise add
    emit(word, merged)
\end{lstlisting}

\subsubsection{Runtime Analysis}
\begin{table}[H]
  \centering
  \begin{tabular}{@{}lcccc@{}}
    \toprule
    \textbf{Approach} & $d=1$ & $d=2$ & $d=3$ & $d=4$ \\
    \midrule
    Stripes (baseline) & 1531 ms & 1543 ms & 2606 ms & 1548 ms \\
    \bottomrule
  \end{tabular}
  \caption{Stripes approach runtime for varying word distance $d$}
\end{table}

\subsubsection{Screenshots}
\begin{figure}[H]
  \centering
  \includegraphics[width=0.9\textwidth]{screenshots/1c_run_1.png}
  \caption{Problem 1(c): Stripes approach -- MapReduce execution start}
\end{figure}

\begin{figure}[H]
  \centering
  \includegraphics[width=0.9\textwidth]{screenshots/1c_run_2.png}
  \caption{Problem 1(c): Stripes approach -- Reduce phase and runtime report}
\end{figure}

\begin{figure}[H]
  \centering
  \includegraphics[width=0.9\textwidth]{screenshots/1c_output.png}
  \caption{Problem 1(c): Stripes co-occurrence matrix output sample}
\end{figure}

% ---------------------------------------------------------------
\subsection{Problem 1(e): Local Aggregation Strategies (5 Marks)}

\subsubsection{Program Logic}
We implemented four specialized Java classes to test two distinct local aggregation strategies across both the Pairs and Stripes architectures.

\textbf{Strategy 1 -- Map-Class Level (Combiner):}
\begin{itemize}
  \item Hadoop's built-in Combiner mechanism is used. The Combiner runs the same logic as the Reducer, but executes \emph{locally on the mapper node} before the shuffle phase.
  \item For Pairs: \texttt{job.setCombinerClass(IntSumReducer.class)} -- trivially reuses the existing reducer.
  \item For Stripes: A custom \texttt{StripesCombiner} class merges \texttt{MapWritable} objects by element-wise addition of counts.
  \item Implementation: \texttt{CoOccurrencePairsCombiner.java}, \texttt{CoOccurrenceStripesCombiner.java}.
\end{itemize}

\textbf{Strategy 2 -- Map-Function Level (In-Mapper Combining):}
\begin{itemize}
  \item A Java \texttt{HashMap} (for Pairs) or \texttt{HashMap<String, MapWritable>} (for Stripes) is created as an instance variable in the Mapper class.
  \item During \texttt{map()}, results are accumulated directly in memory with \textbf{no calls to \texttt{context.write()}}.
  \item Only in the \texttt{cleanup()} method are the fully aggregated results emitted to Hadoop.
  \item This bypasses Hadoop's internal sorting/spilling overhead entirely, making it theoretically the fastest approach.
  \item Implementation: \texttt{CoOccurrencePairsInMapper.java}, \texttt{CoOccurrenceStripesInMapper.java}.
\end{itemize}

\subsubsection{Runtime Comparison}
\begin{table}[H]
  \centering
  \begin{tabular}{@{}lcccc@{}}
    \toprule
    \textbf{Configuration} & $d=1$ & $d=2$ & $d=3$ & $d=4$ \\
    \midrule
    Pairs (no aggregation)          & 1537 ms & 1546 ms & 1554 ms & 1544 ms \\
    Pairs + Combiner                & 1549 ms & 1542 ms & 1557 ms & 1557 ms \\
    Pairs + In-Mapper               & 1535 ms & 1549 ms & 1580 ms & 1544 ms \\
    \midrule
    Stripes (no aggregation)        & 1531 ms & 1543 ms & 2606 ms & 1548 ms \\
    Stripes + Combiner              & 1536 ms & 1547 ms & 1556 ms & 1549 ms \\
    Stripes + In-Mapper             & 1538 ms & 1551 ms & 1557 ms & 1545 ms \\
    \bottomrule
  \end{tabular}
  \caption{Complete runtime comparison across all local aggregation strategies}
\end{table}

\subsubsection{Analysis}
In our local standalone Hadoop environment operating on the relatively small Wikipedia-50-ARTICLES dataset (50 files, $\sim$700 KB total), all configurations converge to approximately the same runtime ($\sim$1540 ms). This is expected because:
\begin{enumerate}
  \item \textbf{No network shuffle overhead}: In local mode, all map and reduce tasks run in the same JVM process. The primary bottleneck that Combiners and In-Mapper Combining address---network I/O during the shuffle phase---does not exist.
  \item \textbf{Small data volume}: With only 50 input files and a filtered vocabulary of 50 words, the number of intermediate key-value pairs is small enough to fit entirely in memory regardless of the aggregation strategy.
  \item \textbf{JVM warm-up dominance}: The $\sim$1.5-second runtime is dominated by Hadoop framework initialization (JVM startup, job setup, class loading) rather than actual data processing.
\end{enumerate}

In a distributed multi-node cluster with large datasets, In-Mapper Combining would be expected to outperform the Combiner, which in turn would outperform the baseline, due to the significant reduction in network transfer during the shuffle phase.

\subsubsection{Screenshots}
\begin{figure}[H]
  \centering
  \includegraphics[width=0.9\textwidth]{screenshots/1e_pc_run.png}
  \caption{Problem 1(e): Pairs + Combiner -- Execution and runtime}
\end{figure}

\begin{figure}[H]
  \centering
  \includegraphics[width=0.9\textwidth]{screenshots/1e_pim_run.png}
  \caption{Problem 1(e): Pairs + In-Mapper Combining -- Execution and runtime}
\end{figure}

\begin{figure}[H]
  \centering
  \includegraphics[width=0.9\textwidth]{screenshots/1e_sc_run.png}
  \caption{Problem 1(e): Stripes + Combiner -- Execution and runtime}
\end{figure}

\begin{figure}[H]
  \centering
  \includegraphics[width=0.9\textwidth]{screenshots/1e_sim_run.png}
  \caption{Problem 1(e): Stripes + In-Mapper Combining -- Execution and runtime}
\end{figure}

% ====================================================================
% PROBLEM 2
% ====================================================================
\newpage
\section{Problem 2: Indexing Documents via Hadoop}

\subsection{Problem 2(a): Document Frequency Calculation (6 Marks)}

\subsubsection{Program Logic}
Implemented in \texttt{DocumentFrequency.java} as a two-job MapReduce chain.

\textbf{Job 1 -- Document Frequency (DF):}
\begin{itemize}
  \item The \texttt{DFMapper} loads stop-words via Distributed Cache and initializes the \textbf{OpenNLP PorterStemmer}.
  \item For each input split (one per document), the mapper tokenizes the text and applies stemming via \texttt{stemmer.stem(token)}.
  \item A \texttt{HashSet<String>} is used to track unique stemmed terms \emph{within each document}. This ensures each term is emitted \textbf{at most once per document}, correctly computing the document frequency.
  \item The reducer sums the counts to get the total DF for each term.
  \item A Combiner (same as reducer) provides local aggregation.
\end{itemize}

\textbf{Job 2 -- Top 100 Extraction:}
\begin{itemize}
  \item Identical \texttt{PriorityQueue} min-heap architecture as Problem 1(a), but with $K=100$.
  \item Output is a TSV file with schema: \texttt{TERM<tab>DF}.
\end{itemize}

\subsubsection{Pseudocode}
\begin{lstlisting}
// Job 1: DFMapper
setup(): load stopwords; stemmer = new PorterStemmer()
map(key, line):
    uniqueTerms = new HashSet()
    for each token in line.split("[^a-z]+"):
        if token not in stopWords:
            stemmed = stemmer.stem(token)
            uniqueTerms.add(stemmed)
    for each term in uniqueTerms:
        emit(term, 1)

// Job 1: DFReducer
reduce(term, counts): emit(term, sum(counts))

// Job 2: Same Top-K logic as Problem 1(a) with K=100
\end{lstlisting}

\subsubsection{Output Sample}
\begin{lstlisting}
TERM            DF
refer           50
s               49
categori        48
apo             46
quot            44
us              43
other           43
includ          42
two             41
link            41
\end{lstlisting}

\subsubsection{Screenshots}
\begin{figure}[H]
  \centering
  \includegraphics[width=0.9\textwidth]{screenshots/2a_run_1.png}
  \caption{Problem 2(a): Document Frequency -- Job 1 (DF extraction) execution}
\end{figure}

\begin{figure}[H]
  \centering
  \includegraphics[width=0.9\textwidth]{screenshots/2a_run_2.png}
  \caption{Problem 2(a): Document Frequency -- Job 2 (Top 100 filtering) and runtimes}
\end{figure}

\begin{figure}[H]
  \centering
  \includegraphics[width=0.9\textwidth]{screenshots/2a_output_1.png}
  \caption{Problem 2(a): Top 100 DF output -- Top portion}
\end{figure}

\begin{figure}[H]
  \centering
  \includegraphics[width=0.9\textwidth]{screenshots/2a_output_2.png}
  \caption{Problem 2(a): Top 100 DF output -- Bottom portion}
\end{figure}

% ---------------------------------------------------------------
\subsection{Problem 2(b): TF-IDF Indexing via Stripes (6 Marks)}

\subsubsection{Program Logic}
Implemented in \texttt{TFIDFIndexing.java} using the Stripes algorithm.

\textbf{Mapper (\texttt{TFMapper}):}
\begin{itemize}
  \item Loads the Top 100 DF TSV file via Distributed Cache (\texttt{Job.addCacheFile()}) into a \texttt{HashMap<String, Integer>} during \texttt{setup()}.
  \item Initializes the PorterStemmer for consistent tokenization.
  \item Extracts the document ID from the input split filename via \texttt{FileSplit.getPath().getName()}.
  \item Builds a \texttt{MapWritable} stripe mapping each Top 100 term found in the document to its term frequency (TF).
  \item Emits \texttt{(DocumentID, \{term $\rightarrow$ TF, ...\})} -- one stripe per split.
\end{itemize}

\textbf{Reducer (\texttt{TFIDFReducer}):}
\begin{itemize}
  \item Also loads the DF cache in \texttt{setup()}.
  \item Merges all incoming stripes for each document by element-wise addition.
  \item For each term present in the merged stripe, computes the TF-IDF score using the formula specified in the assignment:
  \[
    \text{SCORE} = \text{TF} \times \ln\!\left(\frac{10000}{\text{DF}} + 1\right)
  \]
  \item Emits the result in the required TSV schema: \texttt{ID<tab>TERM<tab>SCORE}.
\end{itemize}

\subsubsection{Pseudocode}
\begin{lstlisting}
// Mapper
setup(): load Top100 DF into HashMap; stemmer = new PorterStemmer()
map(key, line):
    docId = inputSplit.filename
    stripe = {}
    for each token in line:
        stemmed = stemmer.stem(token)
        if stemmed in Top100:
            stripe[stemmed] += 1
    if stripe not empty: emit(docId, stripe)

// Reducer
setup(): load Top100 DF into HashMap
reduce(docId, stripes):
    merged = {}
    for each stripe: merged += stripe
    for each (term, tf) in merged:
        df = dfMap[term]
        score = tf * ln(10000/df + 1)
        emit("docId\tterm", score)
\end{lstlisting}

\subsubsection{Output Sample}
\begin{lstlisting}
ID              TERM        SCORE
107301.txt      state       52.015
107301.txt      citi        78.451
108594.txt      region      36.729
108594.txt      citi       126.728
113147.txt      p          241.387
113147.txt      s          207.613
113147.txt      quot      2905.339
1338.txt        american   239.684
1338.txt        state      196.501
\end{lstlisting}

\subsubsection{Screenshots}
\begin{figure}[H]
  \centering
  \includegraphics[width=0.9\textwidth]{screenshots/2b_run_1.png}
  \caption{Problem 2(b): TF-IDF Indexing -- Map phase execution}
\end{figure}

\begin{figure}[H]
  \centering
  \includegraphics[width=0.9\textwidth]{screenshots/2b_run_2.png}
  \caption{Problem 2(b): TF-IDF Indexing -- Reduce phase and runtime}
\end{figure}

\begin{figure}[H]
  \centering
  \includegraphics[width=0.9\textwidth]{screenshots/2b_output_1.png}
  \caption{Problem 2(b): TF-IDF output -- First document scores}
\end{figure}

\begin{figure}[H]
  \centering
  \includegraphics[width=0.9\textwidth]{screenshots/2b_output_2.png}
  \caption{Problem 2(b): TF-IDF output -- Additional document scores}
\end{figure}

% ====================================================================
% ENVIRONMENT
% ====================================================================
\newpage
\section{Environment \& Build Configuration}
\begin{table}[H]
  \centering
  \begin{tabular}{@{}ll@{}}
    \toprule
    \textbf{Component} & \textbf{Version} \\
    \midrule
    Operating System   & macOS (Apple Silicon) \\
    Java (JDK)         & OpenJDK 17 \\
    Apache Hadoop      & 3.3.6 (standalone local mode) \\
    Apache Maven       & 3.9.x \\
    OpenNLP Tools      & 1.9.3 \\
    Dataset            & Wikipedia-50-ARTICLES ($\sim$717 KB) \\
    \bottomrule
  \end{tabular}
  \caption{Development environment specifications}
\end{table}

A comprehensive \texttt{Makefile} was developed to automate compilation and execution of all problems. Example usage:
\begin{lstlisting}[language=bash]
make compile        # Build the JAR via Maven
make 1a             # Run Problem 1(a)
make 1b d=3         # Run Problem 1(b) with distance d=3
make 1c d=4         # Run Problem 1(c) with distance d=4
make 1d_pc d=2      # Run Pairs + Combiner with d=2
make 1d_pim d=2     # Run Pairs + In-Mapper with d=2
make 1d_sc d=2      # Run Stripes + Combiner with d=2
make 1d_sim d=2     # Run Stripes + In-Mapper with d=2
make 2a             # Run Document Frequency extraction
make 2b             # Run TF-IDF Indexing
make clean          # Remove all output directories
\end{lstlisting}

\end{document}
